\subsection{Pregunta 1: análisis del proceso DHCP (DORA)}
\autor{Benjamín}

Se capturó en modo Simulación el intercambio entre el host Wingman y el router Gekko, que actúa como servidor DHCP de la LAN 20.25.16.0/20 (dirección del router: 20.25.16.1). Para forzar una solicitud nueva, el host se pasó a configuración estática y luego otra vez a DHCP.

\evid{Q1_DHCP/dora_simulacion.png}{Mensajes DORA entre Wingman y Gekko en modo Simulación.}

\subsubsection*{Los cuatro mensajes}

Los mensajes viajan en paquetes UDP: el cliente usa el puerto 68 y el servidor el 67. Wingman aún no tiene dirección IP durante todo el intercambio, por lo que todos van con destino de difusión (\texttt{255.255.255.255}) y el switch Gekko~SW los replica a todos sus puertos; Dizzy y Thrash los reciben, pero los descartan porque no son servidores ni los destinatarios de la configuración.

\begin{center}\footnotesize
\setlength{\tabcolsep}{4pt}
\begin{tabular}{@{}>{\raggedright\arraybackslash}p{2.3cm}>{\raggedright\arraybackslash}p{2.3cm}>{\raggedright\arraybackslash}p{1.9cm}>{\raggedright\arraybackslash}p{2.3cm}>{\raggedright\arraybackslash}p{1.6cm}>{\raggedright\arraybackslash}p{3.8cm}@{}}
  \toprule
  Mensaje & Sentido & IP origen & IP destino & Puertos UDP & Función \\
  \midrule
  Discover & Wingman $\rightarrow$ difusión & 0.0.0.0 & 255.255.255.255 & 68 $\rightarrow$ 67 & El cliente busca servidores DHCP en su segmento. \\
  Offer & Gekko $\rightarrow$ Wingman & 20.25.16.1 & 255.255.255.255 & 67 $\rightarrow$ 68 & El servidor propone una configuración de red. \\
  Request & Wingman $\rightarrow$ Gekko & 0.0.0.0 & 255.255.255.255 & 68 $\rightarrow$ 67 & El cliente acepta la oferta; en la capa de aplicación indica el servidor elegido (20.25.16.1). \\
  Acknowledgment & Gekko $\rightarrow$ Wingman & 20.25.16.1 & 255.255.255.255 & 67 $\rightarrow$ 68 & El servidor confirma la asignación. \\
  \bottomrule
\end{tabular}
\end{center}

Como resultado del intercambio, Wingman recibió la dirección \texttt{20.25.16.11}, la máscara \texttt{255.255.240.0} y la puerta de enlace \texttt{20.25.16.1}. La primera dirección asignada es la .11 porque las diez primeras direcciones útiles de la subred (de .1 a .10) están excluidas del pool para asignaciones estáticas.

\pend{confirmar el tipo de cada mensaje en la pestaña Inbound/Outbound PDU Details de Packet Tracer}

\evid{Q1_DHCP/offer_wingman.png}{Offer recibida por Wingman (capas 2, 3, 4 y 7).}
\evid{Q1_DHCP/request_wingman.png}{Wingman recibe la Offer y responde con el Request.}
\evid{Q1_DHCP/request_ack_gekko.png}{Gekko recibe el Request y responde con el Acknowledgment.}

En la simulación, la Offer aparece alrededor de 2 segundos (tiempo simulado) después del Discover.

\subsubsection*{Direcciones IP del Discover}

\evid{Q1_DHCP/discover_capa3.png}{Detalle del DHCP Discover en Gekko: capas 2, 3, 4 y 7.}

En la capa de red, el Discover lleva como \textbf{origen 0.0.0.0} y como \textbf{destino 255.255.255.255}. En la capa de enlace, la MAC de origen es la de Wingman (\texttt{00D0.BC6D.B529}) y la de destino es la de difusión (\texttt{FFFF.FFFF.FFFF}).

Se usan esas direcciones porque el cliente todavía no tiene ninguna dirección IP asignada, así que no puede usar una como origen y usa la dirección ``no especificada'' \texttt{0.0.0.0}. Tampoco conoce la dirección del servidor, por lo que envía el mensaje a la dirección de difusión limitada \texttt{255.255.255.255}, que llega a todos los equipos del segmento sin que un router la reenvíe a otras redes. Así, cualquier servidor DHCP presente en la LAN puede recibirlo y responder.

\subsubsection*{Segmento sin pool DHCP}

Para comprobarlo se eliminó temporalmente el pool en Gekko (\cmd{no ip dhcp pool POOL\_LAN}) y se repitió la solicitud desde Wingman. Packet Tracer informó que la solicitud DHCP falló, y el host quedó con la dirección \textbf{169.254.181.41}. Después se restauró el pool y Wingman volvió a obtener una dirección de la subred 20.25.16.0/20.

Al no haber ningún pool, el router no genera una Offer, por lo que el cliente nunca recibe configuración de red. Tras agotar sus intentos, el host se asigna a sí mismo una dirección del rango \texttt{169.254.0.0/16}, una dirección \textit{link-local} (APIPA) con máscara \texttt{255.255.0.0} \pend{verificar la máscara en la captura}. Este estado significa que la autoconfiguración falló: la dirección solo permite comunicarse con otros equipos del mismo segmento que también se hayan autoasignado una, y no incluye puerta de enlace, por lo que el host no puede salir de su LAN.

\evid{Q1_DHCP/sin_pool.png}{Host solicitando DHCP en un segmento sin pool.}
